\documentclass[11pt]{article}
\usepackage{mathblatt}
% Prüfstein 08.10.2026 abends: Lösungen zum Blatt K4W nach bau/bauregeln.md § 8. Werte nachgerechnet.
\newcommand{\kennung}{K4W}
\begin{document}
\pfheftstil
\fancyfoot[C]{{\footnotesize\color{mbgrau}\kennung\hfill\thepage}}
\pfheftkopf{Hypotenuse berechnen}{Klasse 9 \textperiodcentered{} Lösungen}

\begin{pfloesung}{}
\lz{1a)}{$c$}{}{}
\lz{b)}{$r$}{}{}
\lz{c)}{$s$}{}{}
\lz{d)}{$e$}{}{}
\end{pfloesung}
\begin{pfloesung}{}
\lz{2a)}{9, 16 und 25 Kästchen}{$3^2 = 9$, \ $4^2 = 16$, \ $5^2 = 25$}{}
\lz{b)}{die Quadrate}{$9 + 16 = 25$; \ Seiten: $3 + 4 = 7 \ne 5$}{}
\end{pfloesung}
\begin{pfloesung}{}
\lz{3a)}{$w^2 = u^2 + v^2$}{Hypotenuse $w$}{}
\lz{b)}{$x^2 = y^2 + z^2$}{Hypotenuse $x$}{}
\lz{c)}{$\overline{BC}^2 = \overline{AB}^2 + \overline{AC}^2$}{rechter Winkel bei $A$, Hypotenuse $\overline{BC}$}{}
\lz{d)}{gilt nicht}{kein rechter Winkel}{}
\end{pfloesung}
\begin{pfloesung}{}
\lz{4b)}{$c = 13$ cm}{$c^2 = 5^2 + 12^2 = 169$ \ $\Rightarrow$ \ $c = 13$; \ $12 < 13 < 17$}{}
\lz{c)}{$c = 15$ cm}{$c^2 = 9^2 + 12^2 = 225$ \ $\Rightarrow$ \ $c = 15$}{}
\end{pfloesung}
\begin{pfloesung}{}
\lz{5a)}{$c = \sqrt{97} \approx 9{,}8$ cm}{$c^2 = 4^2 + 9^2 = 97$ \ $\Rightarrow$ \ $c \approx 9{,}849$}{}
\lz{b)}{$c = 130$ cm}{$1{,}2$ m $= 120$ cm; \ $c^2 = 120^2 + 50^2 = 16\,900$}{}
\end{pfloesung}
\begin{pfloesung}{}
\lz{6.}{$d = 51$ cm}{$d^2 = 45^2 + 24^2 = 2\,601$}{}
\end{pfloesung}
\begin{pfloesung}{}
\lz{7.}{$40$ m}{quer: $\sqrt{80^2 + 60^2} = 100$ m; \ $140 - 100 = 40$}{}
\end{pfloesung}
\begin{pfloesung}{}
\lz{8.}{nein}{Diagonale: $\sqrt{40^2 + 30^2} = 50$ cm $< 52$ cm}{}
\end{pfloesung}
\begin{pfloesung}{}
\lz{9.}{$\overline{AB} = \sqrt{1\,300} \approx 36{,}1$ m}{$\overline{AB}^2 = 12^2 + 34^2 = 1\,300$ \ $\Rightarrow$ \ $\overline{AB} \approx 36{,}06$}{}
\end{pfloesung}
\begin{pfloesung}{}
\lz{10.}{$x = \sqrt{29\,156} \approx 170{,}8$ cm}{$x^2 = 170^2 + 16^2 = 29\,156$ \ $\Rightarrow$ \ $x \approx 170{,}75$}{}
\end{pfloesung}
\end{document}
